
\subsubsection{2.1 Execution-Based Code Refinement}

Execution feedback has been used as a signal for code generation and repair. CodeRL \citep{le2022coderl} used unit test execution as a reward signal for reinforcement learning. Self-Debug \citep{Chen2023} demonstrated that LLMs can use execution error traces to iteratively fix their own code. LDB \citep{li2024ldb} showed that runtime debugging traces enable localized bug identification. These methods evaluate execution feedback without comparison against AI-based feedback under matched experimental conditions.

\subsubsection{2.2 AI-Based Feedback Methods}

Self-Refine \citep{madaan2023self} demonstrated that models can iteratively improve outputs through self-generated critique. Constitutional AI \citep{bai2022constitutional} established that AI feedback can guide behavior without human labels. Reflexion \citep{shinn2023reflexion} combined verbal self-reflection with task outcomes. These approaches reduce infrastructure requirements compared to execution sandboxes but may miss errors that only manifest at runtime.

\subsubsection{2.3 Positioning}

Existing work compares complete methods rather than isolating signals. This study contributes an orthogonal comparison: fixing the method and varying only the feedback signal, which reveals that execution's advantage stems from counterfactual localization information.
